\documentclass[letterpaper,12pt]{article}
\usepackage[T1]{fontenc}
\usepackage[utf8]{inputenc}
\usepackage[spanish]{babel}
\usepackage{framed}
\usepackage{amsfonts,setspace}
\usepackage{amsmath}
\usepackage{amssymb, amsmath, amsthm}
\usepackage{comment}
\usepackage[most]{tcolorbox}
\usepackage{amssymb}
\usepackage{dsfont}
\usepackage{anysize}
\usepackage{pdfpages}
\usepackage{multicol}
\usepackage{enumerate}
\usepackage{graphicx}
\usepackage{float}
\usepackage[dvipsnames]{xcolor}
\usepackage[left=1.5cm,top=0.9cm,right=1.5cm, bottom=1.4cm]{geometry}
\usepackage{fancyhdr}
\pagestyle{fancy}
\definecolor{blue(pigment)}{rgb}{0.2, 0.2, 0.6}
% Page
\oddsidemargin -0.4in
\textwidth 7in
\topmargin -0.6in
\textheight 9.1in
\parindent 0em
\parskip 1ex

%Teoremas, Lemas, etc.
\theoremstyle{plain}
\newtheorem{teo}{Teorema}
\newtheorem{lem}{Lemma}
\newtheorem{prop}{Proposici\'on}
\newtheorem{cor}{Corolario}
\theoremstyle{definition}
\newtheorem{defi}{Definici\'on}
\newtheorem{eje}{Ejemplo}
\newtheorem{obs}{Observaci\'on}
% fin macros

% Comandos más usados
% Conjuntos
\newcommand{\N}{\mathbb{N}}
\newcommand{\Z}{\mathbb{Z}}
\newcommand{\Q}{\mathbb{Q}}
\newcommand{\R}{\mathbb{R}}
\newcommand{\C}{\mathbb{C}}
\newcommand{\RR}{\mathbb{R} \times \mathbb{R}}
% Flechas
\newcommand{\ssi}{\Longleftrightarrow}
\newcommand{\imp}{\Longrightarrow}
\newcommand{\pmi}{\Longleftarrow}
% Trigonométricas
\newcommand{\sen}{\text{sen}}

\usepackage{versions}
\includeversion{solution}

\let\epsilon\varepsilon

\begin{document}
\def\Pauta{1}
%%%%%ESCRIBIR 1 si quieren ver la Prueba con Pauta%%%%
%%%%%ESCRIBIR 0 si quieren ver sólo la Ayudantía%%%%
	
	\fancyhead[L]{\itshape{Escuela de Ingeniería}}
	\fancyhead[R]{\itshape{Universidad de O'Higgins}}
	
	\begin{minipage}{11.5 cm}
         \vspace{-6mm}
		\begin{flushleft}
			\vspace{1mm}
			\textbf{ING1002-2 Cálculo Diferencial e Integral}\\
			\textbf{Profesor:}  Gonzalo Flores G.  \\
			\textbf{Ayudante:}  Juan I. Riquelme \& Cristian Acevedo R. \\
		\end{flushleft}
	\end{minipage}
	
	\begin{picture}(2,3)
		\put(430,9.5){\includegraphics[scale=0.22]{Imagenes/UOH.png}}
	\end{picture}
	\vspace{-6mm}
	\begin{center}
		\LARGE \bf{Enunciado Ayudantía 8}
	\end{center}
    \vspace{-8.5mm}
    \begin{center}
		4 de diciembre de 2025
	\end{center}

\begin{enumerate}[\textbf{P\arabic{enumi}.}]
\item las siguientes funciones, encuentre el polinomio de Taylor de orden $n$ en torno al punto $x_0$, $T_f^n(x)$:

\vspace{-0.5cm}
\begin{enumerate}
    \item Taylor de orden 3 de \( f(x) = \dfrac{x^2}{1 - x} \) en torno a \( x = 3 \)

\color{blue}
\textbf{Sol:}
\[
f(x)=\frac{x^2}{1-x},\quad f'(x)=\frac{2x(1-x)+x^2}{(1-x)^2}=\frac{2x-x^2}{(1-x)^2},
\]
\[
f''(x)=\frac{2(1-x)^2-(2x-x^2)2(1-x)(-1)}{(1-x)^4}
=\frac{2(1-x)^2+2(2x-x^2)(1-x)}{(1-x)^4}
=\frac{2(1-x+x^2)}{(1-x)^3},
\]
\[
f'''(x)=\frac{d}{dx}\left(\frac{2(1-x+x^2)}{(1-x)^3}\right)
=\frac{2[-(1-x)^3+3(1-x+x^2)(1-x)^2]}{(1-x)^6}
=\frac{4(2-x)}{(1-x)^4}.
\]
Evaluamos en \(x_0=3\):
\[
f(3)=\frac{9}{1-3}=-\tfrac{9}{2},\quad
f'(3)=\frac{2(3)-9}{(1-3)^2}=\frac{-3}{4},\quad
f''(3)=\frac{2(1-3+9)}{(1-3)^3}=\frac{14}{-8}=-\tfrac{7}{4},
\]
\[
f'''(3)=\frac{4(2-3)}{(1-3)^4}=\frac{-4}{16}=-\tfrac{1}{4}.
\]
Por tanto, el polinomio de Taylor de orden 3 es
\[
T_3(x) = -\tfrac{9}{2} - \tfrac{3}{4}(x-3) - \tfrac{7}{8}(x-3)^2 - \tfrac{1}{24}(x-3)^3.
\]
\color{black}
    
    \item \color{blue}
\textbf{[Desarrollo Propuesto].-} \color{black} Taylor de orden 5 de 
$ f(x) = \ln(1+x) $
centrado en $x = 0$. Encuentre un valor aproximado de $\ln(2)$.
\color{blue}
Este desarrollo es solo para el \textbf{orden $n=2$.}, queda propuesto el de orden $n=5.$
\[
f'(x) = \frac{1}{1+x}, \quad f''(x) = -\frac{1}{(1+x)^2}
\]
Evaluamos en $x=0$:
\[
f(0) = 0, \quad f'(0) = 1, \quad f''(0) = -1
\]
Por tanto, el polinomio de Taylor de orden 2 es:
\[
P_2(x) = f(0) + f'(0)x + \frac{f''(0)}{2}x^2 = x - \frac{x^2}{2}
\]
Para aproximar $\ln(2)$, usamos $x=1$:
\[
\ln(2) \approx P_2(1) = 1 - \frac{1}{2} = \frac{1}{2}
\]
\color{black}
\end{enumerate}



\item Use el método de Newton para estimar el cero de $f(x) = \exp(x) + x + x^2$ con tres iteraciones.
\color{blue}
Quería hacer una aclaración, en la ayudantía les mostré como calcular un cero de $f$ usando el método de newton, que se hace con esta formula:
$$
x_{k+1} = x_k - \frac{f(x_k)}{f'(x_k)}
$$
Desarrollando de esta manera:
\[
f(x) = e^x + x + x^2,
\]
\[
f'(x) = e^x + 1 + 2x.
\]
punto inicial:
\[
x_0 = 0.
\]
\(x_0 = 0\):
\[
f(0) = e^0 + 0 + 0 = 1,
\]
\[
f'(0) = e^0 + 1 + 0 = 2.
\]
\(x_1\):
\[
x_1 = x_0 - \frac{f(0)}{f'(0)}
    = 0 - \frac{1}{2}
    = -0.5.
\]
Evaluamos en \(x_1 = -0.5\):
\[
f(-0.5) = e^{-0.5} - 0.5 + 0.25 \approx 0.35653,
\]
\[
f'(-0.5) = e^{-0.5} + 1 - 1 = e^{-0.5} \approx 0.60653.
\]
Cálculo de \(x_2\):
\[
x_2 = -0.5 - \frac{0.35653}{0.60653}
     \approx -1.08808.
\]
Evaluamos en \(x_2 \approx -1.08808\):
\[
f(-1.08808) = e^{-1.08808} -1.08808 + 1.18393 \approx 0.43263,
\]
\[
f'(-1.08808) = e^{-1.08808} + 1 - 2.17616 \approx -0.83937.
\]

Cálculo de \(x_3\):
\[
x_3 = -1.08808 - \frac{0.43263}{-0.83937}
     \approx -0.57278.
\]
Los valores del método de Newton son:
\[
x_1 = -0.5, \qquad
x_2 = -1.08808, \qquad
x_3 = -0.57278, \quad \ldots
\]
Se observa que no hay convergencia, ya que la función:
\[
f(x) = e^x + x + x^2
\]
satisface:
\[
f(x) > 0 \quad \text{para todo } x \in \mathbb{R},
\]
por lo cual no posee ceros reales, entonces solo nos aproximamos al cero de $f$ a pesar de que nunca lleguemos exactamente a el, esto es un aspecto importante de los métodos númericos, por más que no podamos llegar al punto que queremos, podemos acercarnos a el tanto como queramos con las iteraciones del método de newton, si se cumplen las condiciones apropiadas de $f$.

Volviendo a lo del inicio, en realidad la ayudantía debería decir estimar un cero de $f$ y no el mínimo, pues para el mínimo se usa:
$$
x_{k+1} = x_k - \frac{f'(x_k)}{f''(x_k)}
$$
Para este caso tenemos que utilizando la fórmula de recurrencia:
$$x_{k+1} = x_k - \frac{f'(x_k)}{f''(x_k)}$$ La función es: $f(x) = e^x + x + x^2$, derivamos
    $$f'(x) = e^x + 1 + 2x$$
    $$f''(x) = e^x + 2$$
Dado que $f''(x) > 0$ para todo $x \in \mathbb{R}$, $x_0 = 0$, (recordar que la idea es que este sea cercano al valor que queremos estimar).
$k=0$.
Evaluamos en $x_0 = 0$:
$$f'(0) = e^0 + 1 + 2(0) = 1 + 1 = 2$$
$$f''(0) = e^0 + 2 = 1 + 2 = 3$$
Cálculo de $x_1$:
$$x_1 = x_0 - \frac{f'(x_0)}{f''(x_0)} = 0 - \frac{2}{3}$$
$$x_1 = -\frac{2}{3}$$
Evaluamos en $x_1 = -\frac{2}{3}$:
$$f'\left(-\frac{2}{3}\right) = e^{-2/3} + 1 + 2\left(-\frac{2}{3}\right) = e^{-2/3} - \frac{1}{3}$$
$$f''\left(-\frac{2}{3}\right) = e^{-2/3} + 2$$
$x_2$:
$$x_2 = x_1 - \frac{f'(x_1)}{f''(x_1)} \approx -2/3 - \frac{e^{-2/3} - 1/3}{e^{-2/3} + 2}\approx -0.73832$$
Evaluamos en $x_2 \approx -0.73832$:
$$f'(x_2) \approx e^{-0.73832} + 1 + 2(-0.73832) \approx 0.47805 + 1 - 1.47664 \approx 0.00141$$
$$f''(x_2) \approx e^{-0.73832} + 2 \approx 0.47805 + 2 \approx 2.47805$$
$x_3$:
$$x_3 = x_2 - \frac{f'(x_2)}{f''(x_2)} \approx -0.73832 - \frac{0.00141}{2.47805}$$
$$x_3 \approx -0.73832 - 0.00057$$
$$x_3 \approx -0.73889 \quad \text{luego} \ x_3, \text{es un punto cercano al mínimo de }f.$$
\color{black}

\item Para $f(x) = x^3 - x^2$:
    \begin{itemize}
        \item[a)] Exprese la suma de Riemann de $f$ en el intervalo $[0, 2]$.

        
\color{blue}
\textbf{Sol:} La Suma de Riemann $R_n$ para una partición equiespaciada en $[a, b]$ se define como:
$$
R_n = \sum_{i=1}^{n} f(x_i) \Delta x
$$
\begin{itemize}
    \item Intervalo: $[a, b] = [0, 2]$
    \item Ancho del subintervalo ($\Delta x$):
    $$
    \Delta x = \frac{b-a}{n} = \frac{2-0}{n} = \frac{2}{n}
    $$
    \item Punto de Muestra (Extremo Derecho $x_i$):
    $$
    x_i = a + i \Delta x = 0 + i \left(\frac{2}{n}\right) = \frac{2i}{n}
    $$
\end{itemize}
Evaluamos la función $f(x) = x^3 - x^2$ en $x_i$:
$$
f(x_i) = f\left(\frac{2i}{n}\right) = \left(\frac{2i}{n}\right)^3 - \left(\frac{2i}{n}\right)^2
$$
$$
f(x_i) = \frac{8i^3}{n^3} - \frac{4i^2}{n^2}
$$
Sustituimos $f(x_i)$ y $\Delta x$ en la Suma de Riemann:
$$
R_n = \sum_{i=1}^{n} \left( \frac{8i^3}{n^3} - \frac{4i^2}{n^2} \right) \cdot \frac{2}{n} = \sum_{i=1}^{n} \left( \frac{16i^3}{n^4} - \frac{8i^2}{n^3} \right)
$$
\color{black}
        \item[b)] Calcule el límite $n \to \infty$ de la suma de Riemann.

        
\color{blue}
\textbf{Sol:} $$
\lim_{n \to \infty} R_n = \lim_{n \to \infty} \sum_{i=1}^{n} \left( \frac{16i^3}{n^4} - \frac{8i^2}{n^3} \right)
$$
Separamos las sumas y usamos las fórmulas $\sum i^3 = \frac{n^2(n+1)^2}{4}$ y $\sum i^2 = \frac{n(n+1)(2n+1)}{6}$:
\begin{align*}
\lim_{n \to \infty} R_n &= \lim_{n \to \infty} \left[ \frac{16}{n^4} \sum_{i=1}^{n} i^3 - \frac{8}{n^3} \sum_{i=1}^{n} i^2 \right] \\
&= \lim_{n \to \infty} \left[ \frac{16}{n^4} \cdot \frac{n^2(n+1)^2}{4} - \frac{8}{n^3} \cdot \frac{n(n+1)(2n+1)}{6} \right] \\
&= \lim_{n \to \infty} \left[ 4 \frac{(n+1)^2}{n^2} - \frac{4}{3} \frac{(n+1)(2n+1)}{n^2} \right] \\
&= \lim_{n \to \infty} \frac{4}{n^2}\left[  (n+1)^2 - \frac{1}{3} (n+1)(2n+1) \right] \\
&= \lim_{n \to \infty} \frac{4}{n^2}\left[ \frac{3}{3} (n+1)^2 - \frac{1}{3} (n+1)(2n+1) \right] \\
&= \lim_{n \to \infty} \frac{4}{3n^2}\left[  3(n+1)^2 -(n+1)(2n+1) \right] \\
&= \lim_{n \to \infty} \frac{4}{3n^2}\left[  3n^2+6n+3 -(2n^2+3n+1) \right] \\
&= \lim_{n \to \infty} \frac{4}{3n^2}\left[  n^2-3n+2 \right] \\
&= \lim_{n \to \infty} \frac{4}{3}\left[  \frac{n^2-3n+2}{n^2} \right] \\
&= \lim_{n \to \infty} \frac{4}{3}\left[  1-\frac{3}{n}+\frac{2}{n^2} \right] \\
&= \frac{4}{3} \\
\end{align*}
\textbf{Nota:} En la ayudantía cuando hago el cálculo de 
$$
= \lim_{n \to \infty} \frac{4}{3n^2}\left[  3(n+1)^2 -(n+1)(2n+1) \right]
$$
en la ayudantía habia puesto
$$= \lim_{n \to \infty} \frac{3\cdot4}{n^2}\left[  3(n+1)^2 -(n+1)(2n+1) \right] $$
Mil disculpas por eso, en al final terminaría dando 12 pero eso es incorrecto, pues comparando con otro método:
\[
\int_0^2 (x^3 - x^2)\,dx.
\]
Dado que
\[
\int x^3\,dx = \frac{x^4}{4}, 
\qquad
\int x^2\,dx = \frac{x^3}{3}.
\]
\[
\int_0^2 (x^3 - x^2)\,dx
=
\left[ \frac{x^4}{4} - \frac{x^3}{3} \right]_0^2.
\]
Por TFC
\[
\frac{2^4}{4} - \frac{2^3}{3}
= \frac{16}{4} - \frac{8}{3}
= 4 - \frac{8}{3}
= \frac{12}{3} - \frac{8}{3}
= \frac{4}{3}.
\]
Lo cual es el valor verdadero de la integral que queríamos calcular.
\color{black}



    \end{itemize}
\item Si $x \sen(\pi x) = \displaystyle \int_0^{x^2} f(t)\, dt$, donde $f$ es una función continua, encuentre $f(4)$.


\color{blue}
\textbf{Sol:} Derivamos a ambos lados y aplicamos TFC \begin{align*}
x \sen(\pi x) &= \int_0^{x^2} f(t)\,dt \\
\frac{d}{dx}\left(x \sen(\pi x)\right) &= \frac{d}{dx} \left( \int_0^{x^2} f(t)\,dt \right) \\
\sen(\pi x) + \pi x \cos(\pi x) &= f(x^2) \cdot 2x \\
\frac{\sen(\pi x) + \pi x \cos(\pi x)}{2x} &= f(x^2) \\
f(4) &= \frac{\sen(2\pi) + 2\pi \cos(2\pi)}{2 \cdot 2} = \frac{0 + 2\pi(1)}{4} = \frac{\pi}{2}.
\end{align*}
\textbf{Nota:} la función $f(x^2)$ pasa a ser $f(x)$ pues $x^2$ pasa a ser variable independiente.
\color{black}

\item  Calcule el valor del l\'imite
$$ \lim_{x\to 0} \frac{\displaystyle\int_0^{x^2} \sen(t^2)dt}{1 - e^{x^6}} $$


\color{blue}
\textbf{Sol:} Aplicamos regla de L'Hopital y TFC: \begin{align*}
\lim_{x \to 0} \frac{\int_0^{x^2} \sen(t^2)\,dt}{1 - e^{x^6}} 
&\xrightarrow{L'h} \lim_{x \to 0} \frac{2x \cdot \sen(x^4)}{-6x^5 e^{x^6}} \\
&= \lim_{x \to 0} \frac{2 \sen(x^4)}{-6x^4 e^{x^6}} \\
&\xrightarrow{L'h} \lim_{x \to 0} \frac{8x^3 \cos(x^4)}{-24x^3 e^{x^6} - 36x^9 e^{x^6}} \\
&= \frac{8 \cdot 1}{-24 \cdot 1} = -\frac{1}{3}.
\end{align*}

\color{black}
\newpage
\subsection*{Algunas fórmulas importantes}
\begin{figure}[H]
    \centering
    \includegraphics[width=0.5\linewidth]{Imagenes/tfc.png}
\end{figure}
\begin{figure}[H]
    \centering
    \includegraphics[width=0.7\linewidth]{Imagenes/3a569818-7699-46ef-902b-b0ebb8996a91.jpg}
\end{figure}
\begin{figure}[H]
    \centering
    \includegraphics[width=0.7\linewidth]{Imagenes/91b97700-b93d-4615-b4e9-b2c12b270fb6.jpg}
\end{figure}
\end{enumerate}
\end{document}
\item  Considere la función
    \[
    f(x) = \sen(x^2).
    \]
    Calcule un polinomio de Taylor $T_f$ de orden 3 para $f$ centrado en $\bar{x} = 0$. Pruebe que existe $\xi \in (0, x)$ tal que
    \[
    f(x) = T_f(x) + \frac{f^{(3)}(\xi)x^3}{3!}.
    \]


\item Encuentre un desarrollo limitado para $\sen(x) + \cos(x)$ en torno a 0, cuyo error máximo de aproximación en el intervalo $(-\pi/2, \pi/2)$ sea inferior a $10^{-3}$.

\item Utilice el m\'etodo de Newton para aproximar una ra\'{\i}z de cada una de las siguientes funciones, realizando tres iteraciones
\begin{enumerate}
    \item $f(x) = x^2 - 2$
    \item $f(x) = \cos(x) - x$
\end{enumerate}

\item \begin{enumerate}
    \item Demuestre que entre todos los rectángulos con diagonal dada $d=1$, el que tiene mayor área, es el cuadrado.
    \item Dos lados de un tri\'angulo tienen longitudes $a$ y $b$, y el \'angulo entre ellos es $\theta$ ¿Qu\'e valor de $\theta$ maximizar\'a el \'area del tri\'angulo?
\end{enumerate}

\item \textbf{[Propuesto].-} Una compañía vende máquinas para realizar moños de regalo, a $P$ pesos cada una, la función de demanda de estas máquinas está dada por:
$$
P(q) = 300 - 0,02q
$$
El costo de producción de cada máquina está dado por \$30 cada una y la compañía tiene un costo fijo de \$9000 pesos al año. Hallar la cantidad de máquinas que debe producir la compañía para maximizar las ganancias en un año.
